% SOURCE_AUTHORITY: sga5_fr_workpass.tex, lines 7875--7900 (LF-normalized unit).
% SOURCE_SHA256: 791F4EFFC5E02832D5D77ED03518C8156D6F07E4C8238B03545DB93D883FBB28
\begin{prooftext}[Demonstração]
Como os funtores \(T^n\) comutam com as translações, basta demonstrar a afirmação quando \(X\) é \(J\)-ádico noetheriano. Neste caso, como \(\grs(X)\) é um \(\gr_J(A)\)-módulo graduado noetheriano (5.1.6), deduz-se da hipótese (ii) e de uma variante de um corolário do teorema de SHIH (cf. Apêndice A.3) que, para todo inteiro \(n\), verificam-se as afirmações seguintes:
\[
\text{(a) }\quad T^n(X)\text{ verifica (MLAR),}
\]
\[
\text{(b) }\quad \grs\,T^n(X)\text{ é um }(\gr_J(A),\mathcal D)\text{-módulo graduado noetheriano.}
\]
O lema de Artin-Rees (4.2.6) permite concluir.
\end{prooftext}

\begin{statement}{Corolário 5.3.2}
O \(\delta\)-funtor
\[
T=(T^i)_{i\in\mathbb N}:\Hom(\mathbb N^\circ,\mathcal C)\longrightarrow\Hom(\mathbb N^\circ,\mathcal D)
\]
induz um \(\delta\)-funtor exato \(A\)-linear
\[
\mathcal E_{\mathcal C}\longrightarrow\mathcal E_{\mathcal D},
\]
que comuta com as translações e, portanto, por passagem ao quociente, um novo \(\delta\)-funtor exato \(A\)-linear, denotado também por \((T^i)_{i\in\mathbb N}\),
\[
\mathrm{AR}\text{-}J\text{-adn}(\mathcal C)\longrightarrow
\mathrm{AR}\text{-}J\text{-adn}(\mathcal D).
\]
\end{statement}
